% Приложение: Дорожная карта курса
\chapter{Дорожная карта курса}
\label{app:roadmap}

Дорожная карта показывает, какие темы методички связаны с лекциями и
лабораторными работами, и фиксирует статус наполнения. Заполняется по мере
прохождения курса.

\begin{longtable}{@{}p{0.45\linewidth}p{0.35\linewidth}l@{}}
\toprule
Тема методички & Занятие (лекция / лаб.) & Статус \\
\midrule
\endhead

Доступ к окружению \ref{ch:02-linux-access}         & Лаб.\ 1 (сем.\ 1)      & готово \\
Оболочка и файловая система \ref{ch:03-shell}       & Лаб.\ 1--2, лек.\ 2   & готово \\
Права доступа и пользователи \ref{ch:05-permissions} & Лаб.\ 2 (сем.\ 1)   & готово \\
Справка, конвейеры, пакеты \ref{ch:07-text-pipelines} & Лаб.\ 2--3, лек.\ 2 & готово \\
Редакторы и сессии (vim/tmux/mc) \ref{ch:09-vim}    & Лаб.\ 2, лек.\ 2      & готово \\
Процессы и ресурсы \ref{ch:12-processes}            & Лек.\ 1 (сем.\ 1)    & готово \\
C, сборка, тесты, память, отладка \ref{ch:13-c-compilation} & Лек.\ 1, лаб.\ (C) & готово \\
Docker и изоляция \ref{ch:19-docker-basics}         & Лек.\ 1, лаб.\ (Docker) & готово \\
Git \ref{ch:25-git}                                  & Лаб.\ 4 (сем.\ 1)    & готово \\
Основы безопасности \ref{ch:24-security}            & Лек.\ (сем.\ 2)      & задел \\
Пайплайны / CI \ref{ch:26-ci}                        & Лек.\ 2 (сем.\ 2)    & задел \\
Скриптинг / IDE \ref{ch:27-scripting-ide}           & ---                  & задел \\
\bottomrule
\end{longtable}

% Легенда статусов: «готово» = раздел наполнен; «будет дополнено» = заглушка,
% материал появится с новыми семинарами/лекциями.
